\section{마무리 (30분) — 학술 배경 · 읽을 자료 · Day4 예고}

% =====================================================================
% 마무리 (30분) — 학술 배경 · 읽을 자료 · Day4 예고 — 교재 제6장
% =====================================================================
\begin{frame}{마무리 — 이 교시에서 배우는 것}
\begin{itemize}
\item 오늘 인용한 문헌 아홉 갈래 중 여섯 — "왜 그런가"와 "언제 틀리는가"
\item 읽을 자료 — 필수 4종(과정 전)·심화 13종(30일·90일)·참고 9종
\item 오늘 만든 산출물 ⑪\til{}⑮ 다섯 — 서로를 \textbf{설명하는} 구조 한 장
\item Day2 문서로 되돌아간 것 12건 — 재기준선은 0회
\item Day4 예고 — 전환·종료 ⑯\til{}⑳, 그리고 사후 읽기 과제
\end{itemize}
\keymsg{표준은 "이렇게 하라"까지만 말한다 — 지표는 왜 정보를 잃고, 나쁜 소식은 왜 올라오지 않는가는 문헌이 말한다}
\note{교재 제6장 도입. Day1 제4장·Day2 제6장과 같은 자리다 — 실행·통제 단계에서 학술 문헌이 답하는 두 질문은 지표(왜 SPI는 종료에서 1이 되는가)와 사람(왜 나쁜 소식은 올라오지 않는가)이다. 30분 배분 — 학술 배경 12분, 읽을 자료 5분, 산출물 5종 되짚기·Day2 갱신 6분, Day4 예고·과제 7분. 6교시 실습 ②가 늦게 끝났으면 학술 배경 세 프레임을 각 2분으로 줄이고 Day4 예고는 자르지 않는다 — Day4 첫 시간의 입력이 오늘 문서이기 때문이다. 서지의 [검증 필요] 표기는 과정 리딩리스트(05\_리딩리스트)의 것을 그대로 따랐고, 슬라이드에는 표기를 생략했으니 인쇄 전 원문 대조를 권한다고 말해 둔다. 예상 소요 2분.}
\end{frame}

\begin{frame}{학술 배경 ① — 지표는 왜 정보를 잃는가: Lipke · Christensen}
\begin{tbl}[\scriptsize]{L{2.3cm}L{4.2cm}L{4.6cm}}
문헌 & 발견 & 온담에서 \\ \midrule
Lipke 2003 \textit{Schedule is Different} & SV·SPI는 원가 단위 \ra{} 종료 시 반드시 SV = 0·SPI = 1, 후반 일정 지표는 정보 상실. 해법 = EV를 PV 곡선에 투영해 시간축 ES & ES 8.28개월 · SV(t) $-$0.72 · SPI(t) 0.920 vs SPI 0.943 — 벌어지기 시작한 7월(진척 40\%)이 가장 정보가 많다 \\
Christensen \& Payne 1992 · Christensen \& Heise 1993 & 미 국방 원가정산 계약 표본 — 누적 CPI는 진척 20\% 이후 ±10\% 안정. "CPI는 좋아지지 않는다" & 표본이 다르다(고정가·발주사 관점). 쓸모는 희망 꺾기 — 조정 CPI 0.968이 50\%에서 안정이면 PMB 초과 확정, 문제는 집행 한도 156.0 안인가 \\
\end{tbl}
\keymsg{흔한 오용 둘 — ES를 CPM의 대체로 읽는 것(지연의 위치는 CPM만 안다) / 착시 포함 CPI 1.032를 20\% 규칙으로 외삽하는 것}
\note{§6.1·§6.3. Lipke의 원 논문은 네 쪽짜리 실무 논문(The Measurable News 2003)이고 EVM 40년의 결함 하나를 고쳤다 — 후속으로 Lipke의 Earned Schedule(2009)과 Vanhoucke의 실증이 ES 기반 예측이 SPI 기반보다 일관되게 정확함을 보였다. 세 오독을 짚는다 — ES는 EVM의 대체가 아니라 확장(같은 PV·EV) / ES는 프로젝트 평균이라 지연의 위치를 모른다(온담 $-$15영업일 평균 vs 임계경로 $-$10) / 후반부에만 쓰는 것이 아니라 두 지수가 벌어지는 시점이 가장 정보가 많다. Christensen의 20\% 규칙은 EAC1의 실무적 정당화가 되었지만 표본은 원가정산 계약이며, 온담 같은 고정가·발주사 관점 EVM에 그대로 옮기면 틀린다(§2.5). 규칙의 진짜 쓸모는 "CPI가 곧 회복된다"는 희망을 꺾는 데 있고, 가장 흔한 오용은 착시 포함 CPI(1.032)를 외삽하는 것이다. 예상 소요 4분.}
\end{frame}

\begin{frame}{학술 배경 ② — 완료된 것만 재는 지표의 병리: Vacanti · Reinertsen}
\begin{itemize}
\item Reinertsen 2009 — 이용률 \ra{} 100\%이면 대기 시간 폭증(M/M/1의 $\rho/(1-\rho)$와 Kingman 근사는 다른 식)
\item R-05의 구조 — 발주사 14명 겸임 \til{}100\% \ra{} 결정 요청이 큐에 쌓임. 처방은 "더 열심히"가 아니라 이용률 낮추기
\item Vacanti 2015 — WIP·cycle time·throughput·work item age. Little's Law는 예측이 아니라 \textbf{안정성 진단} 도구
\item flow debt — 오래된 항목 방치 + 새 항목 빨리 끝내기 = 평균 개선, 부채 누적. SPI 수렴과 같은 병리
\item 백분위수(50/85/95\%)로 약속한다 — 평균이 아니라 분포. S10 WIP 41·age 12가 SI EV 하락의 선행 신호
\end{itemize}
\keymsg{처방은 같다 — 완료된 것 옆에 진행 중인 것의 나이를 병기한다(§2.7 흐름 지표, §4.2 결함 나이)}
\note{§6.4. Day2 §6.4가 Reinertsen을 배치 크기·큐·WIP의 경제학으로 소개했다면 Day3에서 그것은 측정이 된다. 큐 논증을 인용할 때 M/M/1과 Kingman 근사를 섞지 않도록 주의시킨다. Vacanti의 세 기여 — Little's Law의 재정의(성립 가정이 깨지는 구간이 불안정의 신호 — 원문 대조 필요), flow debt, Cycle Time Scatterplot과 백분위수. 2교시 §2.7 표(S7\til{}S18)는 이 넷을 EVM 옆에 둔 것이고, 5교시 §4.2가 같은 처방(work item age 병기)을 결함 나이에 적용했다. 학생에게 묻기 — "여러분 팀의 번다운은 완료만 재는가, 진행 중 항목의 나이도 재는가". 예상 소요 3분.}
\end{frame}

\begin{frame}{학술 배경 ③ — 사람: Keil·Smith · Staw · Brooks}
\begin{tbl}[\scriptsize]{L{2.6cm}L{4.1cm}L{4.4cm}}
문헌 & 발견 & 이 교재가 빌린 처방 \\ \midrule
Smith \& Keil 2003 · Keil, Im \& Mähring 2007 (mum effect, 원전 Rosen \& Tesser 1970) & 정보 비대칭 · 비난 예상(blame) · 조직 분위기(climate)가 나쁜 소식의 보고를 가늘게 만든다 & 성격 교정이 아니라 구조 — 트리거를 관측 항목으로, 예외 보고자의 대우로 climate를 정한다(§5.2, R-05 05-12\ra{}07-15) \\
Staw 1976 \textit{Knee-deep in the big muddy} & escalation의 핵심은 sunk cost가 아니라 self-justification — 자기가 내린 결정일수록 추가 투입 & 착수 결정자 $\neq$ 중단 결정자 · 대안은 각각 다른 사람의 결정 · 사전에 정한 숫자(11-13 S22 속도 620)가 자동 발동 \\
Brooks 1975 \textit{Mythical Man-Month} & 늦은 프로젝트에 인력을 더하면 더 늦어진다 — 학습·교육 부담 + 경로 $n(n-1)/2$(82명 = 3,321) & 회복 대안 1은 인력 추가가 아닌 병행·재배치 / 교체는 전임·후임 동시 투입 / 결함 전담 2명은 R3 병행 해제로 \\
\end{tbl}
\keymsg{셋 다 결론은 "사람을 고치지 말고 결정의 구조를 미리 정하라"다 — EX-01의 대안 3개가 그 구조다}
\note{§6.5\til{}6.7. Keil·Smith 계열은 소프트웨어 프로젝트에 mum effect를 이식한 일련의 연구이고, 같은 갈래의 Kappelman, McKeeman \& Zhang(2006) Dominant Dozen이 조기 신호 12개를 사람·프로세스 요인으로 순위화해 상위 신호가 대부분 기술이 아님을 보였다 — §2.8 조기 신호 5개의 이론적 틀. Staw의 실험은 Day1이 게이트의 중단 기준을 착수 결정과 동시에 제출하게 하고 판정자를 분리한 근거이며, Day3에서는 예외 보고의 대안이 "각각 다른 사람의 결정"이어야 하는 이유(§4.6)다 — Flyvbjerg 열 편향의 ⑩. Brooks는 Day2 §3.4 crashing의 한계로 이미 인용했고 오늘 세 자리(회복 계획·인력 교체·결함 전담)에 다시 나타났다 — Williams(2002)의 시스템 다이내믹스가 이 격언을 압축\ra{}병행\ra{}재작업\ra{}추가 압축의 정량 구조로 승격시켰다. Kerzner(§6.8)는 참조서 — 90\% 증후군과 스폰서십 실패 유형만 빌린다. 예상 소요 4분.}
\end{frame}

\begin{frame}{읽을 자료 — 필수(과정 전) · 심화(30일 · 90일)}
\begin{tbl}[\scriptsize]{L{1.2cm}L{4.6cm}L{5.4cm}}
구분 & 자료 & 어느 부분 · 왜 \\ \midrule
필수 & Lipke, \textit{Schedule is Different} (Measurable News, 2003) & 네 쪽 전체 · ES 네 공식을 강의장에서 처음 보면 오후가 무너진다 \\
필수 & ANSI/PMI 19-006-2019 \textit{Standard for EVM} 용어 절 · PRINCE2 7 Progress practice & PMB·BAC·관리예비비 정의 · 허용범위·예외 보고 원문 \\
30일 & Fleming \& Koppelman \textit{EVPM} 4판(PMI 2010) · Christensen 두 논문 · Vacanti 2015 & 고정가 EVM·EAC 전제 · CPI 안정성 표본의 한계 · flow debt·Little's Law \\
30일 & Staw 1976 (\textit{OBHP} 16(1)) · Smith \& Keil 2003 (\textit{ISJ} 13(1)) · Kappelman et al. 2006 & self-justification · 나쁜 소식의 구조 · 조기 신호 12개 \\
90일 & Reinertsen 2009 · SCL Protocol 2판(2017) · PMBOK 8판 M\&C · PRINCE2 7 세 프로세스 · Brooks 『맨먼스 미신』 · Kan 2판 · \textit{Accelerate} + DORA 2025 & 큐·이용률 · 원칙 8·10·14 · 통제 활동 배치 · 발주·수행 PM 분리 · Rayleigh·DRE · 4지표 정의 \\
\end{tbl}
\keymsg{읽지 않아도 되는 것 — 전제 없는 EVM 공식 카드, SPI·CPI 팀 KPI 대시보드 백서, DORA 서열화 사례집, Lewin 3단계}
\note{§6.10. 필수 4종 중 첫째는 사실 Day2 워크북 완성 예시본 5종(60분)이다 — 슬라이드에는 문헌만 실었다. 참고 9종(워크북을 쓰다가 손을 뻗을 것 — ANSI/PMI 19-006 Baseline maintenance, PMI 형상관리 실무표준 2007, PMBOK 6판 7.4.2, PRINCE2 7 Issues·Change practice, Hillson \& Simon ATOM 3판, ISO 21502 §7, 감리기준 고시 2024-53호, SW진흥법·하도급법 조문, SCL 원칙 10)은 교재 표를 가리킨다. "읽지 않아도 되는 것"은 읽은 것이 잘못이 아니라 인용할 때 한계를 함께 말하라는 뜻 — Lewin의 unfreeze–change–refreeze는 Cummings, Bridgman \& Brown(2016)이 Lewin이 만든 적 없음을 사료로 보였다. 감리기준(§6.9)은 민간 온담이 빌린 구조이며 빌릴 때 귀속 판정 원칙을 함께 빌려야 한다는 점만 한 줄 덧붙인다. 예상 소요 4분.}
\end{frame}

\begin{frame}{오늘 만든 다섯 — 성과보고가 중심이고 나머지 넷이 각주다}
\begin{tbl}[\scriptsize]{L{0.5cm}L{3.2cm}L{3.6cm}L{3.6cm}}
\# & 산출물 (워크북 §) & 오늘의 숫자 & ⑫와의 연결 \\ \midrule
⑪ & 변경요청서·영향 분석·CCB 의결서 (§1) & CR-08\til{}11 · +185 FP · +1.02 · CCB 10-15 · BL-02 146.82 & EAC4 ETC 행 "변경 대가 1.02"의 출처 \\
⑫ & 성과보고 EVM·ES 제9호 (§2) & EV 73.43 / AC 71.13 · SPI 0.943 · SPI(t) 0.920 · CPI 1.032\ra{}0.968 · EAC4 150.86 · A10 $-$10 \ra{} EX-01 & 중심 \\
⑬ & 이슈 로그·리스크 등록부 갱신 (§3) & 이슈 11건 · 잔여 EMV 7.26 · 컨틴전시 10.20 = 5.30+2.47+1.41+1.02 & SV 음수의 이유 — IS-02(R-05 실현) \\
⑭ & 품질 통제 기록 (§4) & 결함 1,140/1,024/116 · 밀도 0.132 > 0.13 · 감리 8건 귀속 이의 2 & 8월 SPI 0.913 정체 = 8월 결함 평탄(시험 정지) \\
⑮ & 조달·계약 변경 기록 (§5) & 라이선스 50/50 · 대기 비용 0.42 · 검수 보류 회신 8영업일 · 귀책 7+3+3$-$3 & CV 양수의 이유 — 계약 변경 로그 1번 \\
\end{tbl}
\keymsg{⑫의 어느 숫자가 다른 넷 어디에도 설명되지 않으면 그 숫자는 아직 이해되지 않은 것이다}
\note{워크북 §6 상호 관계 도식과 §0.2 20종 표. Day1의 다섯이 서로를 고쳤고 Day2의 다섯이 서로를 검산했다면 Day3의 다섯은 서로를 설명한다. 반대 방향도 말한다 — ⑫ 없이 나머지 넷은 목록에 그친다(이슈 11건이 임계경로에 무엇을 했는지, 변경 4건이 EAC를 얼마나 움직이는지, 결함 116건이 R3 개발자의 시간을 얼마나 먹는지는 성과보고에서만 한 숫자가 된다). 공통 규칙 셋을 다시 — ① 같은 사건은 같은 날짜·같은 숫자로 다섯 문서에 나타난다(R-05 실현 = IS-02 05-12 · 5월 SI SV $-$2.49 · 합의서 5호 0.42 · 귀책 7일 · S10 velocity 480) ② 착시는 지우지 않고 병기한다 ③ 모든 문서는 결정 요청으로 끝난다. 예상 소요 4분.}
\end{frame}

\begin{frame}{Day2 문서로 되돌아간 것 — 12건, 재기준선 0회}
\begin{tbl}[\scriptsize]{L{3.0cm}L{5.4cm}L{2.6cm}}
Day2 대상 & 갱신 & 결정·시점 \\ \midrule
⑥ 범위·WBS·사전 카드 5장 & v1.1 — FP 12,400 \ra{} 12,585, SI 68.20 \ra{} 69.22, 합계 145.80 \ra{} 146.82 & CCB 10-15 \ra{} BL-02 10-23 \\
⑦ RTM v1.1 & v1.2 — StRS 215 · 변경 이력 열 CR-06\til{}11 · 갱신 주기를 스프린트 리뷰 시로(감리 1번) & P1 PMO, FZ-11 \\
⑧ 일정 기준선 CPM 표 & A12 20\ra{}30일·TF 0, A09 TF 15, 2차 컷오버 03-27 추가. \textbf{A10 $-$10은 지연으로 기록}(EX-01) & P1 PM · 스폰서 10-16 \\
⑨ 원가 기준선 PV·컨틴전시 & 2026-12 +0.33 · 2027-01 +0.69 · 배정표 v1.2 · Rules of Credit 부속 2 정식 첨부 & BL-02 \\
⑩ 등록부 v1.1 · 격자 & v1.2 — 15행 재평가·R-14\til{}R-20 등재·잔여 7.26 / 원가 축 "+10.2 = 컨틴전시 전액"(CR-10) & 10-23, 정미란 서명 \\
\end{tbl}
\keymsg{기준선은 승인된 변경(BL-02)으로만 움직였다 — 지연(A10 $-$10)과 초과(밀도 0.132·클라우드 +0.6)는 편차로 남아 결정을 기다린다}
\note{워크북 §6 갱신 표 12행을 다섯 줄로 압축했다(품질 축 측정 정의·여유 소유권 문안·조정표 되돌아올 조건 해제는 노트로만). Day2 §7이 예고한 그대로 — 오늘의 측정이 어제의 기준선에 갱신을 만들었지만 재기준선(Day2 §3.5 항목 10 — CCB + 이사회, 최대 2회)은 쓰지 않았다. 그림 1-5(편차를 보이게 두기 vs 계획을 고쳐 편차 없애기)를 다시 띄워도 좋다. 학생에게 묻기 — "BL-02에 A10 $-$10을 반영하면 무엇이 달라지는가": 임계경로가 기준선 안으로 들어가 회복 계획의 근거가 사라지고, 다음 달 성과보고는 SPI(t)가 1 근처로 올라가 아무것도 재지 않게 된다. 예상 소요 3분.}
\end{frame}


\begin{frame}{Day4 예고 — 전환·종료: 오픈은 이벤트가 아니라 이관이다}
\begin{tbl}[\scriptsize]{L{0.8cm}L{4.4cm}L{5.6cm}}
 & Day4 산출물 & Day3에서 받는 입력 \\ \midrule
⑯ & 컷오버 런북 · SAC(서비스 수용 기준) & EX-01 결정(1차/2차 창) · 심각도 1·2 미해결 0 · R-14 두 번 컷오버 \\
⑰ & 운영 이관 체크리스트 · SLA/OLA/UC & 이슈 로그 미결 6건 · 감리 대응 완료 여부 · 하도급 5 \\
⑱ & 종료 보고서 · 미해결 사항 이관 목록 & BL-02 대비 최종 EAC · 변경 로그 11건 · 잔여 EMV \\
⑲ & 교훈 등록부 & 조기 신호 5개 · 착시 5개 · 감리 귀속 이의 결과 \\
⑳ & 편익 실현 원장 & 편익 축 G(118억) · B-08 리드타임 12h · 측정 책임자 실명 \\
\end{tbl}
\keymsg{Day4의 첫 질문 — "오늘 종료 보고서를 쓴다면 미해결 사항 이관 목록에 무엇이 들어가는가"}
\note{워크북 §7 "Day4로 이어지는 것" 표. Day4는 컷오버 런북(T$-$24h\til{}T+24h, go/no-go 게이트, point of no return, 롤백 트리거), ITIL Version 5 기준 서비스 전환(SAC·Early Life Support·SLA/OLA/UC 3층·Change Enablement — ITIL 4 어휘로 고정하되 현행은 v5), 편익 실현 원장 6칸(측정 책임자 실명), 종료 보고서와 교훈, 그리고 Product 트랙으로 넘어가는 접합점(이관 목록의 한 줄)을 다룬다. 오늘의 EX-01 결정이 Day4의 컷오버 창을 정하므로, Day4 아침에 "10-16에 스폰서가 무엇을 결정했는가"를 [추가 설정]으로 공개하고 시작한다. 예상 소요 3분.}
\end{frame}

\begin{frame}{사후 읽기 — 30일 안에, 그리고 Day4 전에}
\begin{itemize}
\item Day4 전(필수) — 교재 제6장 6.1 Lipke(2003) ES 원문 요약 4쪽 · 워크북 §7 표 · 자기 성과보고서 1장 재작성(SPI(t) 추가)
\item 30일 — Fleming \& Koppelman 4판 중 PMB·관리예비비 장 · Vacanti 2015 중 work item age 장 · Keil·Smith mum effect 논문 1편
\item 90일 — 자기 프로젝트의 마지막 변경 1건을 6축으로 다시 재고, CCB 의결서가 있으면 네 칸 중 비어 있는 칸을 찾는다
\item 오늘의 산출물 5종(⑪\til{}⑮)은 Day4 아침 09:00에 갤러리에 다시 붙인다 — 종료 보고서의 입력이다
\item 질문은 워크북 각 §X.7 점검 질문에 자기 답을 적어 Day4에 가져온다
\end{itemize}
\keymsg{Day3의 한 문장 — 편차는 측정으로 끝나지 않는다. 결정 요청으로 끝난다}
\note{교재 6.4 읽을 자료 표의 필수·심화(30일)·심화(90일) 구분 그대로. 서지는 리딩리스트 검증 표기를 따른다. "자기 성과보고서 재작성"은 Day4 도입에서 3명에게 발표시킨다(사전 공지). 마지막 문장은 Day1 "숫자 없는 헌장은 소설이다", Day2 "기준선은 발견되지 않는다, 협상된다"와 나란히 벽에 붙는 오늘의 명제다. 예상 소요 3분.}
\end{frame}
